% \definecolor{name}{model}{spec}
\definecolor{freakishgreen}{HTML}{0A982B}
\definecolor{color_bp}{HTML}{3365FF}
\definecolor{color_or}{HTML}{FA2F00}
\definecolor{color_bg}{HTML}{085D04}
\definecolor{color_bc}{HTML}{3365FF}
\definecolor{color_tbl_header}{HTML}{b2dfee} % From: color_aquamarine
\definecolor{color_tbl_row}{HTML}{ECD0F5}

    % \newcommand{\redsout}{\bgroup\markoverwith{\textcolor{red}{\rule[0.5ex]{2pt}{0.4pt}}}\ULon}
    % \newcommand{\orangesout}{\bgroup\markoverwith{\textcolor{orange}{\rule[0.5ex]{2pt}{0.4pt}}}\ULon}
    % \newcommand{\deleteline}[2]{\textsuperscript{\textcolor{black}{ #1}} \textcolor{red}{\sout{#2}}}
    % \newcommand{\addline}[2]{\textsuperscript{\textcolor{black}{ #1}} \textcolor{green}{#2}}

% Temporal, deprecated
% Definición del comando para tablas en formato APA 7
% Argumentos:
% #1: Label (etiqueta)
% #2: Caption (título)
% #3: Column specification (ej: ll|ll)
% #4: Header content (las filas de cabecera con sus \hline)
% #5: Body content (los datos)
% #6: Note (nota al pie de la tabla)
\NewDocumentCommand{\mitabla}{ m m m m m m }{
    \begin{table}[H]
        \centering
        % El entorno threeparttable ajusta las notas al ancho de la tabla automáticamente
        \begin{threeparttable}
            \caption{#2}
            \label{#1}
            \vspace{2mm}
            \rowcolors{3}{color_tbl_row}{} 
            \begin{tabular}{#3}
                \hline
                \rowcolor{color_tbl_header}
                #4 \\
                \hline
                #5
                \hline
            \end{tabular}
            
            \ifx\relax#6\relax\else
                \begin{tablenotes}[flushleft] % Alinea la nota al borde izquierdo de la tabla
                    \small
                    \item \textit{Nota}. #6
                \end{tablenotes}
            \fi
        \end{threeparttable}
    \end{table}
}

% % Comando para tablas APA 7 con posición opcional
% % Argumentos:
% % O{H} : Argumento opcional de posición (Default: H)
% % #2   : Label (etiqueta)
% % #3   : Caption (título)
% % #4   : Column specification (ej: ll|ll)
% % #5   : Header content
% % #6   : Body content
% % #7   : Note (nota al pie de la tabla)
% \NewDocumentCommand{\apaTable}{ O{H} m m m m m m }{
%     \begin{table}[#1]
%         \centering
%         \begin{threeparttable}
%             \caption{#3} % Eliminamos las barras temporales
%             \label{#2}
%             % No hace falta \vspace aquí, el captionsetup hace el trabajo el cual 
%             % se define en la cabecera de main.tex
%             \rowcolors{3}{color_tbl_row}{} 
%             \begin{tabular}{#4}
%                 \hline
%                 \rowcolor{color_tbl_header}
%                 #5 \\
%                 \hline
%                 #6
%                 \hline
%             \end{tabular}       
%             \ifx\relax#7\relax\else
%                 \begin{tablenotes}[flushleft]
%                     \small
%                     \item \textit{Nota}. #7
%                 \end{tablenotes}
%             \fi
%         \end{threeparttable}
%     \end{table}
% }

% \NewDocumentCommand{\apaTable}{ O{htbp} m m m m m m }{
%     \begin{table}[#1]
%         \centering
%         \begin{threeparttable}
%             \caption{#3}
%             \label{#2}
%             % Eliminamos colores para cumplir APA 7
%             \begin{tabular}{#4}
%                 \toprule
%                 #5 \\
%                 \midrule
%                 #6
%                 \bottomrule
%             \end{tabular}       
%             \ifx\relax#7\relax\else
%                 \begin{tablenotes}[flushleft, para]
%                     \small
%                     \textit{Nota}. #7
%                 \end{tablenotes}
%             \fi
%         \end{threeparttable}
%     \end{table}
% }

\NewDocumentCommand{\apaTable}{ O{htbp} m m m m m m }{
    \begin{table}[#1]
        \centering
        \begin{threeparttable}
            \caption{#3}
            \label{#2}
            % Eliminamos colores para cumplir APA 7
            \begin{tabular}{#4}
                \toprule
                #5 \\
                \midrule
                #6
                \bottomrule
            \end{tabular}       
            \ifx\relax#7\relax\else
                \begin{tablenotes}[flushleft, para]
                    \small
                    \textit{Nota}. #7
                \end{tablenotes}
            \fi
        \end{threeparttable}
    \end{table}
}

% test
% Requiere \usepackage{threeparttable} y \usepackage{graphicx}
\NewDocumentCommand{\apaFigure}{ O{htbp} m m m m }{
    \begin{figure}[#1]
        \centering
        % Definimos un ancho fijo (ej: 0.5\textwidth) para alinear todo
        \begin{minipage}{0.5\textwidth} 
            \caption{#3} % Título
            \label{#2}   % Label
            
            \vspace{6pt}
            \includegraphics[width=\textwidth]{#4} % Ruta del PDF
            
            \vspace{6pt}
            \small \textit{Nota}. #5 % Nota
        \end{minipage}
    \end{figure}
}

% Requiere \usepackage{subcaption}
% #1: Label principal
% #2: Título de la figura
% #3: Ancho total del bloque (ej: 0.8\textwidth)
% #4: Contenido de las subfiguras (los entornos \begin{subfigure}...)
% #5: Texto de la nota
\NewDocumentCommand{\apaSubFigures}{ m m m m m }{
    \begin{figure}[htbp]
        \centering
        \begin{minipage}{#3} % Define el ancho donde se alineará el título y la nota
            \caption{#2}
            \label{#1}
            
            \vspace{6pt}
            \centering % Centra las subfiguras dentro de la minipage
            #4 
            
            \vspace{10pt}
            \begin{flushleft}
                \small \textit{Nota}. #5
            \end{flushleft}
        \end{minipage}
    \end{figure}
}
